\newpage
\section{Incidência da empenagem horizontal}

% Item 3 do roteiro: para cada arquivo, achar o it que zera a deflexao de
% profundor na condicao de cruzeiro (mn 0.85, a c 0.5053, d2 pm 0, menu de
% -> variavel 1). Registrar o it final e um print do plano de Trefftz.
% PREENCHER apos as rodadas.

A incidência da empenagem horizontal de cada configuração foi ajustada para
anular a deflexão de profundor na condição de cruzeiro. Em cada arquivo
fixou-se a meta de sustentação do ponto de projeto e a restrição de momento
de arfagem nulo pelo profundor, e a variável de projeto \texttt{it} foi
iterada até a deflexão residual ficar abaixo de $0{,}1^\circ$.

\begin{table}[H]
\centering
\caption{Incidência da empenagem que anula a deflexão de profundor em cruzeiro.}\label{tab:trimagem}
\renewcommand{\arraystretch}{1.25}
\begin{tabular}{lccc}
\toprule
Configuração & $i_t$ [graus] & $\delta_e$ residual [graus] & $\alpha$ [graus] \\
\midrule
CG dianteiro (\texttt{fwd.avl}) & --- & --- & --- \\
CG traseiro (\texttt{aft.avl}) & --- & --- & --- \\
\bottomrule
\end{tabular}
\end{table}

% Inserir os prints do plano de Trefftz dos dois casos:
% \begin{figure}[H] ... images/trefftz_fwd.png / trefftz_aft.png ... \end{figure}
